% Created 2026-09-06 Sun 11:19
% Intended LaTeX compiler: pdflatex
\documentclass[11pt, a4paper, abstract=yes]{scrartcl}
\usepackage[utf8]{inputenc}
\usepackage[T1]{fontenc}
\usepackage{graphicx}
\usepackage{longtable}
\usepackage{wrapfig}
\usepackage{rotating}
\usepackage[normalem]{ulem}
\usepackage{amsmath}
\usepackage{amssymb}
\usepackage{hyperref}
\usepackage{amsthm,amssymb,amsmath}
\usepackage{hyperref}
\usepackage{graphicx}
\theoremstyle{definition}
\newtheorem{definition}{Definition}
\newtheorem{mexample}[definition]{Example}
\theoremstyle{remark}
\newtheorem{remark}[definition]{Remark}
\theoremstyle{plain}
\newtheorem{lemma}[definition]{Lemma}
\newtheorem{claim}[definition]{Claim}
\newtheorem{proposition}[definition]{Proposition}
\newtheorem{theorem}[definition]{Theorem}
\newtheorem{corollary}[definition]{Corollary}
\newtheorem{cor}[definition]{Corollary}
\newtheorem{conjecture}[definition]{Conjecture}
\newtheorem{question}[definition]{Question}
\newcommand{\set}[1]{\left\{{#1}\right\}}
\newcommand\setsuchas[2]{\left\{\,{#1}\,\vrule\,{#2}\,\right\}}
\newcommand{\vek}[1]{\boldsymbol{#1}}
\newcommand{\NN}{{\mathbb{N}}}
\newcommand{\NNp}{{\mathbb{N}_+}}
\newcommand{\ZZ}{{\mathbb{Z}}}
\newcommand{\CC}{{\mathbb{C}}}
\newcommand{\RR}{{\mathbb{R}}}
\newcommand{\QQ}{{\mathbb{Q}}}
\newcommand{\seqnum}[1]{\href{https://oeis.org/#1}{\textrm \underline{#1}}}
\author{Jan Snellman\thanks{Department of Mathematics, Linköping University, 581 83 Linköping, Sweden. \texttt{jan.snellman@liu.se}. ORCID: \href{https://orcid.org/0009-0002-6676-5068}{0009-0002-6676-5068}}}
\usepackage{iftex}
\ifPDFTeX\usepackage[T1]{fontenc}\else\usepackage{fontspec}\fi
\usepackage{orcidlink}
\hypersetup{hidelinks}
\newcommand{\Koz}{\operatorname{Koz}}
\newcommand{\shift}{\operatorname{shift}}
\newcommand{\ff}{\operatorname{ff}}
\newcommand{\inn}{\operatorname{in}}
\newcommand{\conv}{\operatorname{conv}}
\newcommand{\indeg}{\operatorname{indeg}}
\newcommand{\tp}{\operatorname{top}}
\newcommand{\HF}{\operatorname{HF}}
\renewcommand{\vek}[1]{\vec{#1}}
\author{Jan Snellman \orcidlink{0009-0002-6676-5068}}
\date{\today}
\title{Hilbert functions over the exterior algebra and \(f\)-vectors in higher rank\\\medskip
\large II: Generalizing Linusson's count to sumsets of f-vectors}
\hypersetup{
 pdfauthor={Jan Snellman \orcidlink{0009-0002-6676-5068}},
 pdftitle={Hilbert functions over the exterior algebra and \(f\)-vectors in higher rank},
 pdfkeywords={},
 pdfsubject={},
 pdfcreator={Emacs 30.2 (Org mode 9.7.11)}, 
 pdflang={English}}
\begin{document}

\maketitle
\begin{abstract}
Let \(E=\bigwedge(k^n)\) be the exterior algebra on \(n\) generators, let
\(F=\bigoplus_{i=1}^{r}Eg_i\) be a graded free \(E\)-module with \(\deg g_i=d_i\), and let \(M\)
range over the monomial submodules of Amata and Crupi. The Hilbert function of \(F/M\) is a shifted
sum of \(f\)-vectors of simplicial complexes, and a companion paper determines the convex hull of the
set of such sums. Here we count them. For \(r=1\) the count is Linusson's (1999), obtained from a
recursion through compressed complexes; we show that \textbf{no new recursion is needed for general
\(r\)}. Introducing two statistics of an \(f\)-vector --- its top nonvanishing degree and the initial
degree of the corresponding ideal --- we prove that the number of achievable Hilbert functions is a
product of \(r-1\) transfer matrices, one per consecutive generator-degree gap, whose entries are
differences of Linusson's own refined counts \(E^p(n,k)\). The rank therefore enters only as the
length of a matrix product. The bridge is a reformulation of the counting problem, proved from
Amata and Crupi's lex submodules, that replaces a sumset cardinality by a count of tuples subject to
a nearest-neighbour constraint. We give the same theorem for the \emph{proper} variant that the companion
paper's polytope is built from, prove the lemma that makes the restriction legitimate ---
lexification preserves properness --- and derive consequences: a linear recurrence in \(r\) for
constant gaps, polynomiality in \(r\) when all gaps vanish, an exact threshold beyond which the count
degenerates to a plain power, and polynomiality in the number of vertices.
\end{abstract}

\medskip\noindent\textbf{Keywords.} exterior algebra; Hilbert function; \(f\)-vector; simplicial
complex; enumeration; transfer matrix; lex submodule; Kruskal--Katona theorem.

\smallskip\noindent\textbf{2020 Mathematics Subject Classification.} Primary 05A15, 13F55;
Secondary 05E40, 13D40, 06A07.
\medskip

\setcounter{tocdepth}{2}
\tableofcontents
\section{Introduction}
\label{sec:orgcb83101}
\subsection{The rank-one count}
\label{sec:org1d4eb6d}

Let \(k\) be a field and \(E=E_n=\bigwedge(k^n)\) the exterior algebra. Every monomial of \(E\) is
squarefree, so every monomial ideal of \(E\) is the Stanley--Reisner ideal of a downward-closed
family \(\Delta\subseteq2^{[n]}\), and \(\dim_k(E/I_\Delta)_j=f_j(\Delta)\) exactly. Since passing to
an initial ideal preserves Hilbert functions, the set of Hilbert functions of graded quotients of
\(E\) \emph{is} the set of \(f\)-vectors of simplicial complexes on at most \(n\) vertices.

Three classical results describe that set. Kruskal--Katona determines it exactly; Kozlov (1997)
computes its convex hull; and Linusson (1999) \textbf{counts} it. Linusson's method is a recursion on
compressed complexes: peel off the last vertex, split the complex into the part avoiding it and the
link of it, and record how far the first part is filled. The bookkeeping variable of that recursion
--- his \(E^p(n,k)\), the number of \(f\)-vectors that are maximal up to level \(k\) and no further
--- is the object this paper turns out to need.
\subsection{The rank-\(r\) question}
\label{sec:org64dbba5}

Let \(F=\bigoplus_{i=1}^rEg_i\) with \(\deg g_i=d_i\in\mathbb Z\), \(d_1\le\cdots\le d_r\). Amata and
Crupi's monomial submodules are those generated by monomials of \(F\); such an \(M\) is automatically
a direct sum \(I_1g_1\oplus\cdots\oplus I_rg_r\) of monomial ideals of \(E\), and the Hilbert function
of the quotient splits,
\[
  \HF(F/M)(j)\;=\;\sum_{i=1}^{r}f_{j-d_i}(\Delta_i),
\]
a sum of \(r\) independently chosen \(f\)-vectors, each shifted by its own generator degree. Writing
\(\mathcal F_n\) for the set of \(f\)-vectors of downward-closed families on \([n]\) and
\(\shift^d\) for the degree-shift, the achievable set is the sumset
\[
  \mathrm{FF}(\vek d)\;=\;\Bigl\lbrace\sum_{i=1}^{r}\shift^{d_i}(\vek s_i)\;:\;\vek s_i\in\mathcal F_n\Bigr\rbrace .
\]
A companion paper determines \(\conv(\mathrm{FF}(\vek d))\): it is a Minkowski sum of \(r\) shifted copies
of the Kozlov simplex. \textbf{This paper computes \(|\mathrm{FF}(\vek d)|\).}

The problem is a sumset cardinality, and those are usually hard. Two facts show where the difficulty
sits. The map \((\vek s_1,\dots,\vek s_r)\mapsto\sum_i\shift^{d_i}(\vek s_i)\) is injective exactly when the
generator degrees are far enough apart --- precisely, when every gap \(d_{i+1}-d_i\) is at least
\(n+1\) --- and in that regime \(|\mathrm{FF}(\vek d)|=|\mathcal F_n|^r\) and there is nothing to do.
Everything of interest happens when the shifts overlap and distinct tuples collide.
\subsection{The result}
\label{sec:org4b32559}

The answer is a transfer-matrix product, and it is built from Linusson's own refinement rather than
from a new recursion. Attach to \(\vek s\in\mathcal F_n\) two statistics: \(\tp(\vek s)\), the greatest \(j\)
with \(s(j)>0\); and \(\indeg(\vek s)\), the least \(j\) with \(s(j)<\binom nj\), the initial degree
of the corresponding monomial ideal. Let \(A[\tau][\rho]\) count the \(\vek s\in\mathcal F_n\) with
\(\tp(\vek s)=\tau\) and \(\indeg(\vek s)=\rho\). Then (Theorem \ref{org989e5d9})
\[
  \lvert\mathrm{FF}(\vek d)\rvert\;=\;\vek u\cdot T(\delta_2)\,T(\delta_3)\cdots T(\delta_r)\cdot\vek 1,
  \qquad \delta_i=d_i-d_{i-1},
\]
where \(T(\delta)[\tau][\tau']=\sum_{\rho>\tau-\delta}A[\tau'][\rho]\) and \(\vek u\) is the \(\tp\)
marginal of \(A\). And \(A\) is not a new object: Lemma \ref{org3782858} identifies it as a difference of
Linusson's numbers,
\[
  A[\tau][\rho]\;=\;E^n(\tau,\rho-1)-E^n(\tau-1,\rho-1),
\]
so one run of his recursion at \(p=n\) supplies everything. The rank \(r\) enters only as the length
of a matrix product.

The bridge is Theorem \ref{org490b051}, which replaces the sumset by a count of tuples subject to a
constraint relating consecutive entries only:
\[
  \lvert\mathrm{FF}(\vek d)\rvert=\#\bigl\lbrace(\vek s_1,\dots,\vek s_r)\in\mathcal F_n^{\,r}:
   \tp(\vek s_{i-1})<\indeg(\vek s_i)+\delta_i\ \text{ for } i=2,\dots,r\bigr\rbrace .
\]
That is proved from Amata and Crupi's theory of lex submodules, which supplies a canonical
representative for each achievable Hilbert function; the constraint is their Definition 2.7 chaining
condition read off the \(f\)-vectors. Once the constraint is nearest-neighbour, the transfer-matrix
form is a rearrangement.
\subsection{What is proved, and what is only observed}
\label{sec:orgc0101f8}

The distinction is kept explicit throughout, and the reader is asked to take the labels seriously.

\begin{center}
\small
\begin{tabular}{ll}
result & status\\
\hline
Thm. \ref{orga133f2c}, Linusson's recursion & cited; transcription verified against his tables\\
Lem. \ref{org3782858}, the joint distribution & proved; verified entrywise for \(n\le6\)\\
Thm. \ref{org490b051}, the chaining reformulation & proved\\
Thm. \ref{org989e5d9}, the counting theorem & proved; verified, 41 configurations\\
Lem. \ref{orgb9a993b}, lexification preserves properness & proved\\
Thm. \ref{orge76c133}, the proper count & proved; reproduces recorded sequences\\
Prop. \ref{org30e89e0}, recurrence in \(r\) & proved\\
Prop. \ref{org2adde05}, degree bound at \(\delta=0\) & bound proved; exact degrees observed only\\
Thm. \ref{org852a59c}, polynomiality in \(p\) & proved modulo Linusson's Thm. 5.3\\
Obs. \ref{orgd3136f2}, \(\lvert S_n\rvert=F^n(n-2)\) & observed for \(n\le10\); \textbf{not proved}\\
\end{tabular}
\end{center}

The asymptotics in \(n\) --- the analogue of Linusson's own \S2.2 --- are not attempted.
\section{Setup}
\label{sec:orgbfff9fc}
We recall only what is needed; the companion paper has the full development.

\begin{definition}
\(E=E_n=\bigwedge(k^n)=k\langle e_1,\dots,e_n\rangle/(e_ie_j+e_je_i,\ e_i^2)\), with
\(\mathfrak m=(e_1,\dots,e_n)\), and for \(\sigma\subseteq[n]\), \(e_\sigma\) the corresponding
monomial. For \(r\ge1\) and \(d_1\le\cdots\le d_r\) in \(\mathbb Z\), let
\(F=\bigoplus_{i=1}^rEg_i\) with \(\deg g_i=d_i\). A \textbf{monomial submodule} is one generated by
monomials \(e_\sigma g_i\); such an \(M\) decomposes uniquely as \(M=\bigoplus_iI_ig_i\) with each
\(I_i\subseteq E\) a monomial ideal (Amata--Crupi, Definition 2.1). It is \textbf{proper} if no \(I_i\) has
a minimal generator of degree \(\le1\).
\label{orgd425f1c}
\end{definition}

\begin{definition}
A \textbf{downward-closed family} on \([n]\) is \(\Delta\subseteq2^{[n]}\) closed under taking subsets, with
\(f_j(\Delta)=|\lbrace\sigma\in\Delta:|\sigma|=j\rbrace|\). Put
\[
  \mathcal F_n=\lbrace(f_0(\Delta),\dots,f_n(\Delta)):\Delta\ \text{downward closed on }[n]\rbrace,
\]
the void family included and singletons not required, and let
\(S_n=\lbrace \vek s\in\mathcal F_n:s(0)=1,\ s(1)=n\rbrace\) be the \textbf{proper} elements --- the \(f\)-vectors
of simplicial complexes using the full vertex set. For \(d\ge0\), \(\shift^d\) prepends \(d\) zeros.
Finally
\[
  \mathrm{FF}(\vek d)=\Bigl\lbrace\sum_i\shift^{d_i}(\vek s_i):\vek s_i\in\mathcal F_n\Bigr\rbrace,\qquad
  \mathrm{FF}_{\mathrm{pr}}(\vek d)=\Bigl\lbrace\sum_i\shift^{d_i}(\vek s_i):\vek s_i\in S_n\Bigr\rbrace .
\]
\label{orgb41e97d}
\end{definition}
\subsection{Notation: which quantities are tuples}
\label{sec:org8692471}

\textbf{An arrow marks a variable whose value is a tuple; a plain letter is a scalar.} The same convention
is used in the companion article, and it is worth stating in full because two different kinds of
index occur here and the difference matters.

\begin{itemize}
\item \(\vek s\in\mathcal F_n\) is a whole \(f\)-vector; \(\vek d=(d_1,\dots,d_r)\) is the shift vector;
\(\vek u\) and \(\vek 1\) are the row and column vectors of Theorem \ref{org989e5d9}. Each is a tuple, and
each carries an arrow.
\item \textbf{A plain subscript on an arrowed letter selects a member of the \(r\)-tuple}, and that member may
itself be a vector or a scalar: \(\vek s_i\) is the \(i\)-th \(f\)-vector of an \(r\)-tuple, and so
still carries an arrow; \(d_i\) is the \(i\)-th generator degree, a scalar, and does not.
\item \textbf{A parenthesised argument selects a coordinate}, and is always a scalar: \(s(j)\) is the number of
\(j\)-subsets counted by \(\vek s\), \(\HF(F/M)(j)\) is a dimension, \(x(j)\) a single entry of
\(\vek x\). So \(s_i(j)\) --- arrowed letter, tuple subscript, coordinate argument --- is the
\(j\)-th coordinate of the \(i\)-th \(f\)-vector, and is a scalar.
\item Matrices are capitals and never carry an arrow (\(A\), \(T(\delta)\)); their entries are written
\(A[\tau][\rho]\), \(u[\tau]\), also scalars.
\item \textbf{Exception, and it is the usual one:} a subscript on a module or a graded object still means its
homogeneous component --- \(E_j\), \(F_j\), \(I_j\), \(\mathrm{Mon}_j(L)\) are subspaces, not
coordinates of a vector, and follow the standard convention rather than this one.
\end{itemize}

The one place the distinction is load-bearing is Definition \ref{orgb41e97d} against
Theorem \ref{org490b051}: without it, \(s_1\) would mean both "the first coordinate of \(\vek s\)"
(where it equals \(n\) for a proper vector) and "the first \(f\)-vector of the tuple". Under the
convention those are \(s(1)\) and \(\vek s_1\), and cannot be confused.

The dictionary between the algebra and the combinatorics is exact: \(I\mapsto\Delta(I)=\lbrace\sigma:
e_\sigma\notin I\rbrace\) is a bijection from monomial ideals of \(E\) onto downward-closed families,
\(\HF(E/I)\) is the \(f\)-vector of \(\Delta(I)\), and both degenerate cases are included
(\(I=E\mapsto\) the void family, \(I=0\mapsto\) the full simplex). Consequently
\(\mathrm{FF}(\vek d)=\lbrace\HF(F/M):M\ \text{monomial}\rbrace\) and
\(\mathrm{FF}_{\mathrm{pr}}(\vek d)=\lbrace\HF(F/M):M\ \text{proper monomial}\rbrace\).

\begin{remark}
Both sets depend on \(\vek d\) only through the gaps \(\delta_i=d_i-d_{i-1}\ge0\): a global shift
re-indexes \(\mathrm{FF}(\vek d)\) bijectively. Negative generator degrees are therefore allowed
throughout, as they are in Amata and Crupi. What is \textbf{not} optional is that \(\vek d\) be sorted; see
Remark \ref{orgb83d4a5}.
\label{org3777dfb}
\end{remark}
\section{Linusson's recursion}
\label{sec:orgef5314e}
\begin{definition}
\emph{(Linusson 1999, \S5.)} For \(p\ge0\) and \(n\ge-1\), let \(F^p(n)\) be the number of \(f\)-vectors of
simplicial complexes with \(f_1\le p\) and \(f_j=0\) for \(j>n\) --- both degenerate \(f\)-vectors
included --- and let \(E^p(n,k)\) count those among them with \(f_i=\binom pi\) for \(i\le k\) but
\(f_{k+1}<\binom p{k+1}\).
\label{org746bcee}
\end{definition}

\begin{theorem}
\emph{(Linusson 1999, Theorem 5.2.)} For \(0\le k\) and \(1\le n\le p\),
\[
  F^p(n)=\begin{cases}
    1+\sum_{i=0}^{n}E^{p-1}(n,i)\,F^{p-1}(i-1), & n<p,\\
    1+F^p(p-1), & n=p,
  \end{cases}
\]
\[
  E^p(n,k)=\begin{cases}
    \sum_{i=k}^{n}E^{p-1}(n,i)\,E^{p-1}(i-1,k-1), & k\le n<p,\\
    E^p(p-1,k), & k<n=p,\\
    1, & k=n=p,
  \end{cases}
\]
with \(F^p(-1)=1\), \(F^p(0)=2\), \(F^0(n)=2\), \(E^p(n,-1)=1\), \(E^p(n,n)=1\), and
\(E^p(n,k)=0\) for \(k>n\). Moreover \(F^p(n)=\sum_kE^p(n,k)\).
\label{orga133f2c}
\end{theorem}

\begin{remark}
The transcription above was checked entry by entry against Linusson's own appendix: all \(100\)
printed entries of his Table 5 and all \(99\) of his Table 6 are reproduced exactly, and an
independent brute-force enumeration of downward-closed families agrees for every \(F^p(n)\) and
\(E^p(n,k)\) with \(p\le5\). One wrinkle is worth recording. At \(k=n=p\) the clause
"\(f_{k+1}<\binom p{k+1}\)" read literally is vacuously false, since \(\binom p{p+1}=0\); but
Linusson's stated boundary value \(E^p(n,n)=1\) and his Table 6 diagonal both require the clause to
be dropped there. We follow Linusson.
\label{org1660dbd}
\end{remark}

Two consequences we use repeatedly. Setting \(p=n\) gives \(F^n(n)=|\mathcal F_n|\); and the
proper count is a difference of two such.

\begin{remark}
\(|S_n|=F^n(n)-F^{n-1}(n-1)\), giving
\[
  \lvert S_n\rvert = 1,\,2,\,5,\,16,\,70,\,457,\,4908,\,95248,\,3617645,\,286007155
  \qquad(n=1,\dots,10),
\]
whose first five terms agree with direct enumeration. The sequence \(F^p(p)\) is
OEIS A007695. A curious identity, \textbf{observed for \(n\le10\) and not proved}: \(|S_n|=F^n(n-2)\); for
instance \(|S_9|=F^9(7)=3617645\), a value printed in Linusson's Table 5. Equivalently, via
\(F^n(n)=1+F^n(n-1)\), it says \(F^n(n-1)-F^n(n-2)=F^{n-1}(n-1)-1\). We do not use it.
\label{orgd3136f2}
\end{remark}
\section{Two statistics, and the joint distribution}
\label{sec:orgb4dc9d1}
\begin{definition}
For \(\vek s\in\mathcal F_n\) put
\[
  \indeg(\vek s)=\min\lbrace j:s(j)<\tbinom nj\rbrace,\qquad
  \tp(\vek s)=\max\lbrace j:s(j)>0\rbrace,
\]
with \(\indeg(\vek s)=n+1\) for the full simplex and \(\tp(\vek s)=-1\) for the void family. Let
\[
  A[\tau][\rho]=\#\lbrace \vek s\in\mathcal F_n:\tp(\vek s)=\tau,\ \indeg(\vek s)=\rho\rbrace,
  \qquad \tau=-1,\dots,n,\quad \rho=0,\dots,n+1 .
\]
\label{org33aa35f}
\end{definition}

Note \(\indeg(\vek s)\le\tp(\vek s)+1\) always, since \(s(\tp(\vek s)+1)=0<\binom n{\tp(\vek s)+1}\) whenever the
is positive; so \(A\) is supported on \(\rho\le\tau+1\).

\begin{lemma}
With the convention \(E^n(-2,\cdot)=0\),
\[
  A[\tau][\rho]\;=\;E^n(\tau,\rho-1)-E^n(\tau-1,\rho-1).
\]
\label{org3782858}
\end{lemma}

\begin{proof}
Work at \(p=n\). Every \(\vek s\in\mathcal F_n\) has \(f_1\le n\); conversely every \(f\)-vector counted by
\(E^n(m,k)\) has \(f_1\le n\), hence belongs to a complex on at most \(n\) vertices, which embeds in
\([n]\). So the objects counted by \(E^n(m,k)\) are exactly certain elements of \(\mathcal F_n\).
Under that identification Linusson's clause "\(f_j=0\) for \(j>m\)" says precisely \(\tp(\vek s)\le m\),
and his clause "\(f_i=\binom ni\) for \(i\le k\) but \(f_{k+1}<\binom n{k+1}\)" --- with the boundary
convention of Remark \ref{org1660dbd} --- says precisely \(\indeg(\vek s)=k+1\). Hence
\[
  E^n(m,k)=\#\lbrace \vek s\in\mathcal F_n:\tp(\vek s)\le m,\ \indeg(\vek s)=k+1\rbrace ,
\]
and subtracting the count at \(m=\tau-1\) from the count at \(m=\tau\) leaves \(\tp(\vek s)=\tau\)
exactly. For \(\tau=-1\) the subtrahend would count \(f\)-vectors with \(\tp\le-2\), of which there
are none --- even the void \(f\)-vector has \(\tp=-1\) --- which is the stated convention.
\end{proof}

So one run of Theorem \ref{orga133f2c} at \(p=n\) delivers the entire joint distribution. Its
\(\indeg\) marginal telescopes to \(E^n(n,\rho-1)\), which is Linusson's Table 6.

\begin{remark}
Lemma \ref{org3782858} was also checked entrywise, for \(n\le6\) and with no mismatches, against
direct enumeration of the downward-closed families themselves. For the record, \(A\) at \(n=6\), rows \(\tau=-1,\dots,6\) and
columns \(\rho=0,\dots,7\):
\label{orge9c0315}
\end{remark}

\[
\begin{pmatrix}
1&0&0&0&0&0&0&0\\
0&1&0&0&0&0&0&0\\
0&5&1&0&0&0&0&0\\
0&20&14&1&0&0&0&0\\
0&43&84&19&1&0&0&0\\
0&25&143&84&14&1&0&0\\
0&1&25&43&20&5&1&0\\
0&0&0&0&0&0&0&1
\end{pmatrix}
\]
\section{The chaining reformulation}
\label{sec:org3b10bdb}
This section replaces a sumset cardinality by a count of tuples with a nearest-neighbour constraint.
The mechanism is that Amata and Crupi's lex submodules give a canonical representative for each
achievable Hilbert function, and that being a lex submodule is a condition on consecutive components
only.

\begin{lemma}
Let \(I\subseteq E\) be a monomial ideal and \(\vek s=\HF(E/I)\in\mathcal F_n\).
\begin{enumerate}
\item For every integer \(t\ge0\): \(\mathfrak m^t\subseteq I\) if and only if \(\tp(\vek s)<t\).
\item \(\indeg(I)=\indeg(\vek s)\).
\end{enumerate}
\label{orgd47ba5e}
\end{lemma}

\begin{proof}
(1) \(\mathfrak m^t\) is spanned by the monomials of degree \(\ge t\), so
\(\mathfrak m^t\subseteq I\) iff \(I_j=E_j\) for all \(j\ge t\), iff \(s(j)=0\) for all \(j\ge t\),
iff \(\tp(\vek s)<t\). This covers \(t=0\) (then \(\mathfrak m^0=E\subseteq I\) iff \(I=E\) iff
\(\tp(\vek s)=-1<0\)) and \(t>n\) (then \(\mathfrak m^t=0\) and \(\tp(\vek s)\le n<t\), both sides true). For
\(t<0\) the equivalence fails, the left side being vacuously true and the right side impossible ---
which is exactly why sortedness of \(\vek d\) is needed in Theorem \ref{org490b051}.
(2) \(\dim I_j=\binom nj-s(j)\), so the least \(j\) with \(I_j\ne0\) is the least \(j\) with
\(s(j)<\binom nj\). For \(I=0\) both sides are \(n+1\) by convention; any value exceeding \(n\) gives
the same truth value in (1), so the convention is harmless.
\end{proof}

\begin{theorem}
Fix \(n\ge1\), \(r\ge1\) and integers \(d_1\le d_2\le\cdots\le d_r\), with
\(\delta_i=d_i-d_{i-1}\). Then
\[
  \lvert\mathrm{FF}(\vek d)\rvert=\#\bigl\lbrace(\vek s_1,\dots,\vek s_r)\in\mathcal F_n^{\,r}\;:\;
  \tp(\vek s_{i-1})<\indeg(\vek s_i)+\delta_i\ \text{ for } i=2,\dots,r\bigr\rbrace .
\]
\label{org490b051}
\end{theorem}

\begin{proof}
(a) By Definition \ref{orgd425f1c} and the dictionary, \(\mathrm{FF}(\vek d)\) is the set of \(\HF(F/M)\)
for \(M\) a monomial submodule. Indeed, for \(M=\bigoplus_iI_ig_i\),
\[
  \HF(F/M)(j)=\sum_i\dim(E/I_i)_{j-d_i}=\sum_i\shift^{d_i}(\vek s_i)(j) ,
\]
and as \((I_1,\dots,I_r)\) ranges over all tuples of monomial ideals --- unconstrained by one
another --- \((\vek s_1,\dots,\vek s_r)\) ranges over all of \(\mathcal F_n^{\,r}\).

(b) Restricted to \textbf{lex} ideals, \(I\mapsto\HF(E/I)\) is a bijection onto \(\mathcal F_n\). It is well
defined (a lex ideal is monomial); injective, since a lex ideal is degreewise the lex-initial segment
of \(\mathrm{Mon}(E_j)\) of size \(\binom nj-s(j)\) and is therefore determined by \(\vek s\); and
surjective by Kruskal--Katona in its exterior form, every \(f\)-vector being attained by a compressed
family, whose complementary ideal is lex.

(c) By Amata and Crupi's Definition 2.6 \(\Leftrightarrow\) 2.7, a lex submodule of \(F\) is exactly
an \(L=\bigoplus_iL_ig_i\) with every \(L_i\) a lex ideal and
\(\mathfrak m^{\rho_i+d_i-d_{i-1}}\subseteq L_{i-1}\) for \(i=2,\dots,r\), where
\(\rho_i=\indeg(L_i)\). The exponent is \(\ge0\) because \(\rho_i\ge0\) and \(d_i\ge d_{i-1}\), so
Lemma \ref{orgd47ba5e} applies and turns the side condition into
\(\tp(\vek s_{i-1})<\indeg(\vek s_i)+\delta_i\). With (b) componentwise, \(L\mapsto(\vek s_1,\dots,\vek s_r)\) is a
bijection from lex submodules onto the chained tuples.

(d) \(L\mapsto\HF(F/L)\) is a bijection from lex submodules onto \(\mathrm{FF}(\vek d)\). Injective: by
Definition 2.6, \(\mathrm{Mon}_j(L)\) is the initial segment of \(\mathrm{Mon}_j(F)\) of size
\(\dim F_j-\HF(F/L)(j)\) in the module order, and an initial segment is determined by its
cardinality, so \(L\) is determined degreewise. Surjective: monomial submodules are graded, so
\(\mathrm{FF}(\vek d)\subseteq\lbrace\HF(F/M):M\ \text{graded}\rbrace=\lbrace\HF(F/L):L\ \text{lex}\rbrace\)
by Amata and Crupi's Theorem 4.2, implication (a) \(\Leftrightarrow\) (c); and conversely a lex
submodule is monomial, so the reverse inclusion holds by (a).

Composing (c) and (d) gives the claim.
\end{proof}

\begin{remark}
Three hypotheses deserve to be stated rather than left implicit.

\textbf{The order convention.} Amata and Crupi's module order makes \(g_1\)'s monomials the largest, so a
lex submodule fills \(g_1\) first and the chaining condition constrains the \textbf{earlier} component
\(I_{i-1}\) using the \textbf{later} component's initial degree. The inequality in Theorem \ref{org490b051}
is oriented accordingly.

\textbf{Sortedness.} \(d_1\le\cdots\le d_r\) is required as stated, because it is what makes every exponent
\(\rho_i+\delta_i\) non-negative, which is the range in which Lemma \ref{orgd47ba5e}(1) holds. This
costs nothing: \(\mathrm{FF}(\vek d)\) is symmetric in its summands, so an unsorted shift vector should be
sorted first. The chained count as written is simply not the right-hand side for unsorted \(\vek d\).

\textbf{Which half of Theorem 4.2.} Only implication (a) \(\Leftrightarrow\) (c) is used --- existence and
uniqueness of a lex representative. That half is sound. The clause (b) of the same theorem is
\textbf{false as printed} for tied initial generator degrees, a defect treated in the third paper of this
series; nothing here depends on it.
\label{orgb83d4a5}
\end{remark}
\section{The counting theorem}
\label{sec:orgb7feaea}
\begin{theorem}
Fix \(n\ge1\), \(r\ge1\) and integers \(d_1\le\cdots\le d_r\) with gaps \(\delta_i=d_i-d_{i-1}\).
Index vectors and matrices by \(\tau,\tau'\in\lbrace-1,0,\dots,n\rbrace\), and let
\[
  u[\tau]=\sum_\rho A[\tau][\rho],\qquad
  T(\delta)[\tau][\tau']=\sum_{\rho>\tau-\delta}A[\tau'][\rho],
\]
with \(A\) as in Definition \ref{org33aa35f} --- equivalently, by Lemma \ref{org3782858}, assembled from one
run of Linusson's recursion at \(p=n\). Then
\[
  \lvert\mathrm{FF}(\vek d)\rvert\;=\;\vek u\cdot T(\delta_2)\cdot T(\delta_3)\cdots T(\delta_r)\cdot\vek 1 .
\]
\label{org989e5d9}
\end{theorem}

\begin{proof}
By Theorem \ref{org490b051}, \(|\mathrm{FF}(\vek d)|\) counts tuples with
\(\tp(\vek s_{i-1})<\indeg(\vek s_i)+\delta_i\). Constraint \(i\) involves \(s_{i-1}\) only through \(\tp\) and
\(\vek s_i\) only through \(\tp\) and \(\indeg\). Summing over tuples in the order \(\vek s_1,\vek s_2,\dots,\vek s_r\)
and grouping each \(\vek s_i\) by its \((\tp,\indeg)\) class therefore gives the stated product: \(u[\tau]\)
counts the admissible \(\vek s_1\) with \(\tp=\tau\), and \(T(\delta_i)[\tau][\tau']\) counts the \(\vek s_i\)
with \(\tp=\tau'\) whose \(\indeg=\rho\) satisfies \(\tau<\rho+\delta_i\).
\end{proof}

The rank enters only as the length of the product. One computation of \(A\) at fixed \(n\) answers
every \(r\) and every shift vector at once, which is what makes values far outside brute-force range
available: \(|\mathrm{FF}(8;0,0)|=27655902\) is immediate, while direct enumeration of the sumset
becomes infeasible around \(n=6\), \(r=2\).

\begin{corollary}
For \(r=1\), \(|\mathrm{FF}|=\sum_\tau u[\tau]=F^n(n)=|\mathcal F_n|\) --- Linusson's own count.
\label{org1eb04b0}
\end{corollary}

\begin{proposition}
\(|\mathrm{FF}(\vek d)|=|\mathcal F_n|^{\,r}\) if and only if every gap satisfies \(\delta_i\ge n+1\).
At \(\delta_i=n\) the constraint fails for exactly one pair, namely
\((\vek s_{i-1},\vek s_i)=(\text{full simplex},\text{void family})\).
\label{org4ddf98c}
\end{proposition}

\begin{proof}
The constraint \(\tau<\rho+\delta\) is vacuous exactly when \(\delta>\max\tau-\min\rho\). Over
\(\mathcal F_n\), \(\max\tp=n\) (full simplex) and \(\min\indeg=0\) (void family), so the threshold is
\(\delta\ge n+1\); at \(\delta=n\) the single violating pair is the one realising both extremes.
\end{proof}

For \(r=2\) this gives \(|\mathcal F_n|^2-1\) at \(\delta=n\), i.e. \(99\), \(675\), \(9215\) for
\(n=3,4,5\), which matches the computed table.

\begin{proposition}
If all gaps vanish then \(T(0)\) is unipotent, and \(|\mathrm{FF}(n;0,\dots,0)|\) is a polynomial in
\(r\) of degree at most \(n+1\).
\label{org2adde05}
\end{proposition}

\begin{proof}
\(T(0)[\tau][\tau']\ne0\) requires some \(\rho>\tau\) with \(A[\tau'][\rho]>0\); since
\(\rho\le\tau'+1\) always, this forces \(\tau'\ge\tau\), so \(T(0)\) is triangular in the
\(\tp\)-ordering. Its diagonal entries are \(T(0)[\tau][\tau]=A[\tau][\tau+1]=1\), the unique
\(s\) with \(\tp=\tau\) and \(\indeg=\tau+1\) being the full truncation at \(\tau\). A unipotent
matrix has \(T^{r-1}\) polynomial in \(r\) of degree at most the matrix size minus one.
\end{proof}

\begin{remark}
The \textbf{exact} degrees are observed rather than proved. Finite differences for \(n=2,\dots,7\) give
degree exactly \(n+1\) in the improper case --- so the bound of Proposition \ref{org2adde05} is tight ---
and exactly \(n-1\) in the proper case of Section \ref{sec:orgb1f50e3}. We state these as observations. For
\(n=2\) the improper values are \(14,30,55,91\) for \(r=2,\dots,5\), a cubic; the proper values are
\(|\mathrm{FF}_{\mathrm{pr}}(2;0^r)|=r+1\).
\label{orgee21856}
\end{remark}
\section{The proper variant}
\label{sec:orgb1f50e3}
The companion paper's polytope is built from \(S_n\), not \(\mathcal F_n\): its Hilbert functions come
from \textbf{proper} monomial submodules. The count restricts accordingly, but the restriction is not free.
The argument of Theorem \ref{org490b051} goes through the lex representative, so what must be checked
is that the lex representative of a properly achievable Hilbert function is itself proper.

\begin{lemma}
Let \(\vek x=\sum_i\shift^{d_i}(\vek s_i)\) with every \(\vek s_i\in S_n\). Then the unique lex submodule
\(L\) with \(\HF(F/L)=\vek x\) is proper.
\label{orgb9a993b}
\end{lemma}

\begin{proof}
By Definition 2.6, \(\mathrm{Mon}_j(L)\) is the initial segment of \(\mathrm{Mon}_j(F)\) of size
\(D_j=\dim F_j-x(j)\), and the segment fills the slots \(g_1,g_2,\dots\) in order. So \(L\) is proper
exactly when, for every \(i\), the segment does not reach \(g_i\)'s slot in the two module degrees
\(j=d_i\) and \(j=d_i+1\) --- that is, when \(D_j\le\sum_{k<i}\binom n{j-d_k}\) there. Writing
\(\binom nm=0\) and \(s_k(m)=0\) outside \(0\le m\le n\),
\[
  D_j=\sum_k\Bigl(\tbinom n{j-d_k}-s_k(j-d_k)\Bigr).
\]
For \(k\ge i\) and \(j\le d_i+1\) we have \(j-d_k\le d_i+1-d_k\le1\), and a proper \(\vek s_k\) has
\(s_k(0)=1=\binom n0\) and \(s_k(1)=n=\binom n1\), while both vanish in negative degrees; so every
such summand is zero. Hence \(D_j=\sum_{k<i}\bigl(\binom n{j-d_k}-s_k(j-d_k)\bigr)\le
\sum_{k<i}\binom n{j-d_k}\), which is exactly the capacity of the slots preceding \(g_i\).
\end{proof}

\begin{theorem}
Let \(A'\) be \(A\) with the columns \(\rho=0\) and \(\rho=1\) set to zero, \(\vek u'\) its row sums, and
\(T'(\delta)\) built from \(A'\) as in Theorem \ref{org989e5d9}. Then
\[
  \lvert\mathrm{FF}_{\mathrm{pr}}(\vek d)\rvert
  =\#\bigl\lbrace(\vek s_1,\dots,\vek s_r)\in S_n^{\,r}:\tp(\vek s_{i-1})<\indeg(\vek s_i)+\delta_i\bigr\rbrace
  = \vek u'\cdot T'(\delta_2)\cdots T'(\delta_r)\cdot\vek 1 .
\]
\label{orge76c133}
\end{theorem}

\begin{proof}
\(S_n=\lbrace \vek s\in\mathcal F_n:\indeg(\vek s)\ge2\rbrace\), which is what deleting the two columns
implements. As in Theorem \ref{org490b051}(d), \(L\mapsto\HF(F/L)\) is injective on lex submodules; it
carries proper lex submodules into \(\mathrm{FF}_{\mathrm{pr}}(\vek d)\), since a proper lex submodule is a
proper monomial submodule, and onto it: given \(\vek x\in\mathrm{FF}_{\mathrm{pr}}(\vek d)\), Theorem 4.2
(a) \(\Leftrightarrow\) (c) supplies the lex representative and Lemma \ref{orgb9a993b} shows it is
proper. Definition 2.6 \(\Leftrightarrow\) 2.7 with Lemma \ref{orgd47ba5e}, restricted componentwise
to lex ideals of initial degree \(\ge2\), identifies proper lex submodules with \(S_n\)-chained
tuples. The factorisation is then the rearrangement of Theorem \ref{org989e5d9} verbatim.
\end{proof}

\begin{proposition}
\(|\mathrm{FF}_{\mathrm{pr}}(\vek d)|=|S_n|^{\,r}\) if and only if every gap satisfies
\(\delta_i\ge n-1\).
\label{org7f42707}
\end{proposition}

\begin{proof}
As in Proposition \ref{org4ddf98c}, with \(\max\tp=n\) and \(\min\indeg=2\) over \(S_n\).
\end{proof}

\begin{remark}
Two sequences recorded in earlier work on this project are reproduced exactly by
Theorem \ref{orge76c133} and extended. For \(r=2\), \(n=2,\dots,9\):
\[
  \delta=0:\quad 3,\ 12,\ 69,\ 627,\ 9904,\ 296289,\ 18009389,\ 2346218172,
\]
\[
  \delta=1:\quad 4,\ 22,\ 162,\ 1858,\ 38232,\ 1571774,\ 138881518,\ 27635025278 .
\]
The prefixes through \(n=6\) are the independently enumerated ones. Neither sequence, nor \(|S_n|\)
itself, is currently in the OEIS.
\label{org937f09b}
\end{remark}
\section{Growth}
\label{sec:org62303f3}
\begin{proposition}
For a constant gap \(\delta\), \(a_r:=|\mathrm{FF}(n;0,\delta,2\delta,\dots)|=\vek u\cdot
T(\delta)^{r-1}\cdot\vek 1\) satisfies the linear recurrence whose characteristic polynomial is
that of \(T(\delta)\), and \(a_r\) grows like \(C\lambda^r\) with \(\lambda\) the Perron root of
\(T(\delta)\). The same holds for the proper count with \(T'(\delta)\).
\label{org30e89e0}
\end{proposition}

\begin{proof}
Cayley--Hamilton applied to \(T(\delta)\), the vectors \(\vek u\) and \(\vek 1\) being fixed;
Perron--Frobenius for the growth rate, \(T(\delta)\) being non-negative.
\end{proof}

\begin{mexample}
At \(n=2\), \(\delta=1\), the characteristic polynomial of \(T(1)\) is
\(\lambda(\lambda-1)(\lambda^2-4\lambda+1)\), so \(a_r=5a_{r-1}-5a_{r-2}+a_{r-3}\); the values are
\(5,20,76,285,1065\) for \(r=1,\dots,5\), and indeed \(5\cdot285-5\cdot76+20=1065\). The Perron root
is \(2+\sqrt3\approx3.732\). At \(n=3\), \(\delta=1\), the polynomial is
\(\lambda^4-9\lambda^3+21\lambda^2-15\lambda+3\), giving
\(a_r=9a_{r-1}-21a_{r-2}+15a_{r-3}-3a_{r-4}\), checked against \(10,69,426,2532,14847\).
\label{orgb96524e}
\end{mexample}

\begin{remark}
Two structural facts about the Perron root, both verified numerically for \(n\le6\) and stated as
observations. As \(\delta\to n+1\) the improper root tends to \(|\mathcal F_n|\), which is
Proposition \ref{org4ddf98c} in the limit; and the \textbf{proper} Perron root equals \(|S_n|\) exactly for
every \(\delta\ge n-1\), which is Proposition \ref{org7f42707}.
\label{org587fcb9}
\end{remark}

The remaining growth question is in the number of vertices. The count \(|\mathrm{FF}(\vek d)|\) itself
cannot be polynomial in \(n\), since it dominates \(|\mathcal F_n|=F^n(n)\). The right generalisation
is Linusson's own: keep the dimension bound and free the vertex bound.

\begin{definition}
For \(p\ge n\), let \(\mathcal F_n^{(p)}\) be the set of \(f\)-vectors with \(f_1\le p\) and
\(f_j=0\) for \(j>n\), and \(\mathrm{FF}^{(p)}(n;\vek d)\) the corresponding sumset.
\label{orgbf5507f}
\end{definition}

\begin{theorem}
With \(A^{(p)}[\tau][\rho]=E^p(\tau,\rho-1)-E^p(\tau-1,\rho-1)\) and \(\vek u^{(p)},T^{(p)}\) built from it
as before,
\[
  \lvert\mathrm{FF}^{(p)}(n;\vek d)\rvert\;=\;\vek u^{(p)}\cdot T^{(p)}(\delta_2)\cdots T^{(p)}(\delta_r)\cdot\vek 1 ,
\]
and for fixed \(n\), \(r\) and \(\vek d\) this is a polynomial in \(p\).
\label{org852a59c}
\end{theorem}

\begin{proof}
Lemma \ref{org3782858} and Theorems \ref{org490b051} and \ref{org989e5d9} go through verbatim over the exterior
algebra on \(p\) generators, with one extra ingredient: the lex representative must respect the
truncation, i.e. \(\mathfrak m^{n+1}F\subseteq M\) must force \(L\supseteq\mathfrak m^{n+1}F\). This
is the mirror image of the capacity argument of Lemma \ref{orgb9a993b}. In module degree
\(j=d_i+m\) with \(m\ge n+1\), every component \(k\le i\) has \(s_k(j-d_k)=0\) because
\(\tp(s_k)\le n\); hence \(D_j\ge\sum_{k\le i}\binom p{j-d_k}\) and the initial segment covers
\(g_i\)'s slot. Polynomiality in \(p\) is then immediate from Theorem \ref{org989e5d9}'s formula, each
entry \(E^p(\tau,\rho-1)\) being a polynomial in \(p\) for fixed \(\tau,\rho\) by Linusson's
Theorem 5.3.
\end{proof}

\textbf{Status.} The reduction above is proved; the polynomiality of the entries is Linusson's Theorem 5.3,
which we cite rather than reprove, so Theorem \ref{org852a59c} is proved modulo that citation. It was
verified against brute-force truncated sumsets for \(n\le3\), \(p\le5\), \(r\le3\), with no
mismatches, and the empirical degrees in \(p\) at \(r=2\) are \(3,5,6,9\) for
\((n,\delta)=(2,0),(2,1),(3,0),(3,1)\).

\begin{remark}
The asymptotics in \(n\) --- transporting Linusson's Theorem 5.3 degree formula and Proposition 5.4
leading coefficient into a leading-coefficient statement for \(|\mathrm{FF}^{(p)}|\) --- are bounded
but real work and are \textbf{not attempted here}.
\label{org1f988a8}
\end{remark}
\section{Verification, tables, and open problems}
\label{sec:orgeaa5788}
\subsection{How the results were checked}
\label{sec:org52dc8b6}

Every claim above labelled "verified" was checked against a route independent of the formula being
tested. Specifically:

\begin{itemize}
\item \textbf{Against brute force.} Direct enumeration of downward-closed families on \([n]\), forming the
sumset explicitly, with no recursion and no Kruskal--Katona: agreement in all \(41\)
configuration-and-variant pairs with \(n\le6\), \(r\le4\) and gaps \(0\) to \(6\), including
asymmetric gap vectors such as \((0,1,3)\) and \((0,2,4)\), for both the proper and the improper
count.
\item \textbf{Against Linusson's own tables.} All \(199\) printed entries of his Tables 5 and 6, exactly
(Remark \ref{org1660dbd}).
\item \textbf{Against a computer-algebra oracle.} Amata and Crupi's \texttt{ExteriorModules} package for Macaulay2
enumerates lex submodules directly. Its counts \(43\), \(69\), \(99\) and \(236\) for \(n=3\) with
\(d=(0,0)\), \((0,1)\), \((-1,2)\) and \((0,0,1)\) are reproduced --- note that \((-1,2)\) exercises
the negative-degree hypothesis --- as is \(368\) for \(n=4\), \(d=(0,1)\).
\item \textbf{Against previously recorded values.} The two sequences of Remark \ref{org937f09b}.
\item \textbf{Internal consistency.} The \(r=1\) specialisation (Corollary \ref{org1eb04b0}) and the product-regime
identities (Propositions \ref{org4ddf98c} and \ref{org7f42707}) for \(n=2,\dots,6\),
\(r\le4\), both variants.
\end{itemize}

No mismatch was found in any of these.
\subsection{Open problems}
\label{sec:orgff6f035}

\begin{enumerate}
\item \textbf{Asymptotics in \(n\)}, as in Remark \ref{org1f988a8}.
\item \textbf{The exact degrees at \(\delta=0\).} Proposition \ref{org2adde05} proves the bound \(n+1\);
Remark \ref{orgee21856} observes that it is attained, and that the proper degree is exactly
\(n-1\). Both want proofs.
\item \textbf{The identity \(|S_n|=F^n(n-2)\)} of Remark \ref{orgd3136f2}, observed for \(n\le10\).
\item \textbf{A bijective proof.} Theorem \ref{org490b051} is proved by way of a canonical representative;
a direct bijection between \(\mathrm{FF}(\vek d)\) and the chained tuples, not routed through lex
submodules, would explain the constraint combinatorially rather than algebraically.
\item \textbf{Beyond constant gaps.} The formula handles an arbitrary sorted \(\vek d\), but the growth results of
Section \ref{sec:org62303f3} assume a constant gap. Non-constant gap vectors give products of distinct
matrices, for which no closed form has been sought.
\end{enumerate}
\subsection{Integer sequences}
\label{sec:orgb2286af}

Six sequences arising here appear not to be in the OEIS: \(|S_n|\) itself, and the \(r=2\)
proper and improper counts at \(\delta=0\) and \(\delta=1\), together with one \(r=3\) example. Of
these \(|S_n|=1,2,5,16,70,457,4908,95248,3617645,286007155\) --- the number of \(f\)-vectors of
simplicial complexes on exactly \(n\) vertices --- seems the most natural candidate for submission,
being the first-difference sequence of A007695, which has no separate entry. Nothing has been
submitted.
\section*{Acknowledgements}
\label{sec:org38cb087}
Svante Linusson's 1999 paper is not merely cited here but used structurally: the refinement he
introduced for his own recursion turns out to be exactly the statistic the rank-\(r\) count requires,
and the present paper would have no content without it. Aliaksandra Kupreyeva's Bachelor's thesis
(Linköping University, 2019; report LiTH-MAT-EX--2019/05--SE), written under my supervision, is where
counting questions about \(f\)-vectors entered this project.

Just as Geheimrat Goethe had the diligent Friedrich Wilhelm Riemer to help him write \emph{Die
Wahlverwandtschaften}, just as Grothendieck had Jean Dieudonné to help him complete EGA, I have
\textbf{Claude}! What exactly that amounted to is set out in the disclosure section below.

Computations were carried out in SageMath, in Python, and in Macaulay2, the latter using Amata and
Crupi's own \texttt{ExteriorModules} package.
\section*{Disclosure of AI assistance}
\label{sec:org89ecfb8}
This manuscript was produced with substantial assistance from a large language model, Claude
(Anthropic), and the extent of it is set out here rather than left to be inferred. It is more
extensive in this paper than in its companion.

The model: identified that Linusson's refined count \(E^p(n,k)\) is the joint-distribution object the
rank-\(r\) problem needs, and proved Lemma \ref{org3782858}; found and proved the chaining reformulation
of Theorem \ref{org490b051}, including the dictionary of Lemma \ref{orgd47ba5e} and the hypotheses
recorded in Remark \ref{orgb83d4a5}; assembled and proved Theorem \ref{org989e5d9} and its consequences
(Corollary \ref{org1eb04b0}, Propositions \ref{org4ddf98c} and \ref{org2adde05}); found and proved
Lemma \ref{orgb9a993b}, without which the proper variant of Theorem \ref{orge76c133} would not follow;
derived the growth results of Section \ref{sec:org62303f3}, including Theorem \ref{org852a59c}; transcribed
Linusson's Theorem 5.2 from the printed page and validated it against both of his appendix tables;
wrote all the supporting code; and drafted the exposition. The author set the problem, directed the
scope, and required the separation of proved from verified results that Section \ref{sec:orgcb83101}'s table
records.

Two corrections the model made to its own earlier work are worth recording, as evidence of what the
verification was for: an initial value of \(|S_9|\) was wrong by one, caught by checking against
OEIS A007695 and against Linusson's printed table; and a threshold in Propositions \ref{org4ddf98c}
and \ref{org7f42707} had been stated for the wrong one of the two variants.

Every counting claim in this paper was checked along a route independent of the formula being tested,
as set out in Section \ref{sec:orgeaa5788}, and the reference list has been checked against primary
sources. The author has verified the results to the best of his ability, made the final edits, and
assumes full responsibility for any errors, mistakes or gaps that remain.
\section*{Funding}
\label{sec:orga83c28e}
This research received no external funding. It was carried out as part of the author's regular duties
at Linköping University, within the fraction of that employment allocated to research.
\section*{Conflicts of interest}
\label{sec:org95a1c21}
The author declares no financial or non-financial conflicts of interest. The nature and extent of
large-language-model assistance in producing this manuscript is disclosed in full in the section
above.
\section*{References}
\label{sec:org6385e97}
\begin{enumerate}
\item L. Amata, M. Crupi, \emph{Hilbert functions of graded modules over an exterior algebra: an algorithmic
approach}, Int. Electron. J. Algebra \textbf{27} (2020), 271--287. DOI 10.24330/ieja.663116.
\item L. Amata, M. Crupi, \emph{A generalization of Kruskal--Katona's theorem}, An. Şt. Univ. Ovidius
Constanţa *28*(2) (2020), 35--51. DOI 10.2478/auom-2020-0018.
\item L. Amata, M. Crupi, \emph{ExteriorModules: a package for computing monomial modules over an exterior
algebra}, J. Softw. Algebra Geom. \textbf{11} (2021), 71--81. DOI 10.2140/jsag.2021.11.71.
\item A. Aramova, J. Herzog, T. Hibi, \emph{Gotzmann theorems for exterior algebras and combinatorics},
J. Algebra \textbf{191} (1997), 174--211. DOI 10.1006/jabr.1996.6908.
\item G. F. Clements, B. Lindström, \emph{A generalization of a combinatorial theorem of Macaulay},
J. Combin. Theory \textbf{7} (1969), 230--238.
\item G. O. H. Katona, \emph{A theorem of finite sets}, in: Theory of Graphs, Academic Press, 1968,
209--214. [Page range as recorded in this project's \texttt{references.bib}; the Tihany proceedings were
co-published by Akadémiai Kiadó and Academic Press and the commonly quoted range is 187--207.
*To be resolved against the primary source before submission.*]
\item D. N. Kozlov, \emph{Convex hulls of \(f\)- and \(\beta\)-vectors}, Discrete Comput. Geom. \textbf{18} (1997),
421--431. DOI 10.1007/PL00009326.
\item J. B. Kruskal, \emph{The number of simplices in a complex}, in: Mathematical Optimization Techniques,
Univ. of California Press, 1963, 251--278.
\item A. Kupreyeva, \emph{Simplices of \(f\)-vectors}, Bachelor's thesis LiTH-MAT-EX--2019/05--SE,
Linköping University, 2019.
\item S. Linusson, \emph{The number of \(M\)-sequences and \(f\)-vectors}, Combinatorica \textbf{19} (1999),
255--266.
\item OEIS Foundation Inc., \emph{The On-Line Encyclopedia of Integer Sequences}, sequence A007695.
Published electronically at \texttt{oeis.org}.
\item J. Snellman, \emph{Hilbert functions of graded modules over an exterior algebra, and Minkowski
sums of Kozlov simplices}, in preparation. [Companion paper; article 1 of this series.]
\item R. P. Stanley, \emph{Combinatorics and Commutative Algebra}, 2nd ed., Progr. Math. 41,
Birkhäuser, 1996.
\end{enumerate}
\end{document}
